\documentclass[12pt,a4paper]{article}

\usepackage{fontawesome5}
\providecommand{\faPencilAlt}{\faIcon{pencil-alt}}

\usepackage{styles/couleurs}
\usepackage{styles/mise_en_page}
\usepackage{styles/boites}
\usepackage{styles/maths}
\usepackage{styles/figures}

\usepackage{enumitem}
\usepackage{array}
\usepackage{tabularx}
\usepackage{amsmath}
\usepackage{amssymb}
\usepackage{graphicx}
\usepackage{tikz}
\usepackage{xparse}

\usetikzlibrary{calc,angles,quotes,arrows.meta}

\setlist[itemize]{leftmargin=1.4em,itemsep=0.15em,topsep=0.2em}
\setlist[enumerate]{leftmargin=*,itemsep=0.22em,topsep=0.25em}

\newcommand{\lignepointillee}{\noindent\dotfill\par\vspace{0.45em}}

\newcommand{\lignesreponse}[1]{%
\begin{tcolorbox}[
enhanced,colback=white,colframe=NeutreC!40,boxrule=0.5pt,
arc=1.5mm,left=3mm,right=3mm,top=2mm,bottom=1mm,
before skip=2pt,after skip=2pt]
\foreach \i in {1,...,#1}{\lignepointillee}
\end{tcolorbox}}

\newcommand{\titreseq}[1]{%
\begin{center}
{\sffamily\large\bfseries\color{SavoirC}#1}
\end{center}}

\begin{document}

\titreseq{Activité 1 -- Ce que nous savons déjà... et ce qui manque}

\section*{Partie A -- Un calcul que nous savons faire}

Dans le triangle \(GHI\) rectangle en \(H\), on connaît :

\[
GH=6\text{ cm}
\qquad\text{et}\qquad
HI=8\text{ cm}.
\]

Calcule \(GI\). Indique la propriété utilisée.

\lignesreponse{9}

\section*{Partie B -- Le problème de la rampe}

Une rampe rectiligne mesure \(5\) m.
Elle forme un angle de \(28^\circ\) avec le sol horizontal.

On cherche la longueur au sol \(AB\).

\begin{center}
\begin{tikzpicture}[scale=1]
  \coordinate (A) at (0,0);
  \coordinate (B) at (6,0);
  \coordinate (C) at (6,3.2);
  \draw[line width=1pt] (A)--(B)--(C)--cycle;
  \draw (5.65,0)--(5.65,0.35)--(6,0.35);
  \draw pic[draw=ActionC,fill=ActionC!20,angle radius=9mm,"$28^\circ$"] {angle=B--A--C};
  \node[below left] at (A) {$A$};
  \node[below right] at (B) {$B$};
  \node[above right] at (C) {$C$};
  \node[above left] at ($(A)!0.5!(C)$) {$5$ m};
  \node[below] at ($(A)!0.5!(B)$) {$?$};
\end{tikzpicture}
\end{center}

\begin{enumerate}
  \item Le théorème de Pythagore permet-il de calculer \(AB\) avec ces seules données ?
  \item Quelles informations connaît-on dans le triangle ?
  \item Quelle longueur cherche-t-on par rapport à l'angle de \(28^\circ\) ?
  \item Propose une estimation raisonnable de \(AB\), sans chercher une valeur précise.
\end{enumerate}

\lignesreponse{4}

\lignesreponse{7}

\section*{Partie C -- Un cas particulier accessible}

Construis un triangle équilatéral de côté \(10\) cm puis trace une hauteur.

Tu obtiens deux triangles rectangles possédant un angle de \(60^\circ\).

\begin{enumerate}
  \item Quelle est la longueur de l'hypoténuse de l'un de ces triangles ?
  \item Quelle est la longueur du côté adjacent à l'angle de \(60^\circ\) ?
  \item Par quel nombre faut-il multiplier l'hypoténuse pour obtenir ce côté adjacent ?
\end{enumerate}

\lignesreponse{16}

\aretenir{Avec une longueur et un angle aigu, il manque un outil nouveau.
Dans le cas particulier de \(60^\circ\), le côté adjacent vaut la moitié de l'hypoténuse.}

\end{document}
